%%%%%%%% COMPUTER VISION 842 - ASSIGNMENT 04 %%%%%%%%%%%%%%%%%

\documentclass{article}

\usepackage{microtype}
\usepackage{graphicx}
\usepackage{booktabs} % for professional tables
\usepackage{hyperref}
\usepackage{subcaption}
\usepackage{caption}


% Attempt to make hyperref and algorithmic work together better:
\newcommand{\theHalgorithm}{\arabic{algorithm}}

% Use the following line for the initial blind version submitted for review:
% \usepackage{icml2025}

% If accepted, instead use the following line for the camera-ready submission:
\usepackage[accepted]{icml2025}

% For theorems and such
\usepackage{amsmath}
\usepackage{amssymb}
\usepackage{mathtools}
\usepackage{amsthm}

% if you use cleveref..
\usepackage[capitalize,noabbrev]{cleveref}

%%%%%%%%%%%%%%%%%%%%%%%%%%%%%%%%
% THEOREMS
%%%%%%%%%%%%%%%%%%%%%%%%%%%%%%%%
\theoremstyle{plain}
\newtheorem{theorem}{Theorem}[section]
\newtheorem{proposition}[theorem]{Proposition}
\newtheorem{lemma}[theorem]{Lemma}
\newtheorem{corollary}[theorem]{Corollary}
\theoremstyle{definition}
\newtheorem{definition}[theorem]{Definition}
\newtheorem{assumption}[theorem]{Assumption}
\theoremstyle{remark}
\newtheorem{remark}[theorem]{Remark}

% Todonotes is useful during development; simply uncomment the next line
%    and comment out the line below the next line to turn off comments
%\usepackage[disable,textsize=tiny]{todonotes}
\usepackage[textsize=tiny]{todonotes}


% The \icmltitle you define below is probably too long as a header.
% Therefore, a short form for the running title is supplied here:
\icmltitlerunning{CV842 - ASSIGNMENT 04}

\begin{document}

\onecolumn

\icmltitle{Computer Vision 2026: Assignment 04}

\begin{icmlauthorlist}
\icmlauthor{Michael Bezuidenhout}{yyy}
\end{icmlauthorlist}

\section{Image Retrieval for Image Re-Identification}

\subsection{Approach}

A small image retrieval system was trained on the 
VeRi dataset \cite{XinchenLiu16a} for vehicle re-
identification. The aim was to train a model that is capable of 
producing an embedding space where images of similar vehicles are clustered 
relatively close to each other (and relatively far away from images of notably 
different vehicles). During evaluation, query and gallery images are embedded 
and each query image is matched to gallery images in order of proximity within 
the embedding space. A capable model should be able to rank gallery images of 
the particular vehicle above vehicles of similar colour and shape (not to 
mention vehicles that are completely different).

\subsection{Model Architecture}

\subsection{Triplet Mining Strategy}

\subsection{Dataset}

\subsection{Training}

\begin{figure}[H]
    \centering
    \includegraphics[width=0.5\textwidth]{images/alpha02-losscurve.pdf}
    \caption{Semi-hard triplet loss curve for training with $\alpha = 0.2$.}
\end{figure}

\subsection{Results and Discussion}

\begin{figure}[H]
    \centering
    \includegraphics[width=0.75\textwidth]{images/alpha02-rankings.pdf}
    \caption{Semi-hard triplet loss curve for training with $\alpha = 0.2$.}
\end{figure}

\subsection{Recommendations}

\nocite{XinchenLiu16}
\nocite{XinchenLiu16a}
\nocite{XinchenLiu18}
\nocite{He15}

\bibliography{main}
\bibliographystyle{icml2025}

%%%%%%%%%%%%%%%%%%%%%%%%%%%%%%%%%%%%%%%%%%%%%%%%%%%%%%%%%%%%%%%%%%%%%%%%%%%%%%%
%%%%%%%%%%%%%%%%%%%%%%%%%%%%%%%%%%%%%%%%%%%%%%%%%%%%%%%%%%%%%%%%%%%%%%%%%%%%%%%
% APPENDIX
%%%%%%%%%%%%%%%%%%%%%%%%%%%%%%%%%%%%%%%%%%%%%%%%%%%%%%%%%%%%%%%%%%%%%%%%%%%%%%%
%%%%%%%%%%%%%%%%%%%%%%%%%%%%%%%%%%%%%%%%%%%%%%%%%%%%%%%%%%%%%%%%%%%%%%%%%%%%%%%

%%%%%%%%%%%%%%%%%%%%%%%%%%%%%%%%%%%%%%%%%%%%%%%%%%%%%%%%%%%%%%%%%%%%%%%%%%%%%%%
%%%%%%%%%%%%%%%%%%%%%%%%%%%%%%%%%%%%%%%%%%%%%%%%%%%%%%%%%%%%%%%%%%%%%%%%%%%%%%%


\end{document}

